\documentclass[11pt]{article}
\usepackage[margin=1in]{geometry}
\usepackage{amsmath,amssymb,amsthm,mathtools}
\usepackage{graphicx}
\usepackage{microtype}
\usepackage[colorlinks=true,linkcolor=blue,citecolor=blue,urlcolor=blue]{hyperref}

\newtheorem{theorem}{Theorem}
\newtheorem{lemma}{Lemma}
\newcommand{\Sn}{\mathbb S^{n-1}}
\newcommand{\cA}{\mathcal A}
\newcommand{\cH}{\mathcal H}

\title{A non-even near-ball branch for the\\
$L_p$--Christoffel--Minkowski equation}
\author{Candidate partial result for arXiv:2411.17896}
\date{13 August 2026}

\begin{document}
\maketitle

\noindent\fbox{\parbox{0.94\linewidth}{\textbf{Status.}
Candidate partial result, likely valid.  The global non-even uniqueness
question is \emph{not} claimed.  We prove a full non-even local branch,
uniqueness already under $C^0$ closeness to the ball, and automatic symmetry
of that branch for even data.}}

\section{Source question}

For a support function $h$ on $\Sn$, write
\[
  \cA h=\nabla^2_{\Sn}h+hI
\]
and let $s_j(h,\cdot)$ be the normalized $j$th elementary symmetric
function of the eigenvalues of $\cA h$.  Patsalos and Saroglou study
\begin{equation}\label{eq:PDE}
  h^{1-p}s_j(h,\cdot)=g,
  \qquad 1\le j\le n-2,\quad 0\le p<1. 
\end{equation}
Their Theorem 1.4 proves small-data uniqueness in the even class and asks
whether evenness can be removed \cite{PS}.

\begin{figure}[h]
  \centering
  \includegraphics[width=0.98\linewidth]{figures/open_question_crop.png}
  \caption{Theorem 1.4 and the non-evenness question on page 5 of
  arXiv:2411.17896.}
\end{figure}

The literal ``natural question'' about a universal intrinsic-volume
$L_p$--Brunn--Minkowski inequality, appearing earlier in the source, is
answered negatively by the source's own Theorem 1.1.  It is not our target.

\section{Proof intuition}

Evenness is not needed at the linearized level.  At the ball, the only new
modes introduced by dropping evenness are the degree-one spherical
harmonics, corresponding to translations.  Their linearized eigenvalue is
$1-p$, which is strictly positive.  Constants also have positive eigenvalue,
while all degrees at least two have negative eigenvalue.  Thus the full
nonlinear map is locally invertible on all H\"older functions.

There is a second point.  A solution assumed only $C^0$-close to the ball is
automatically H\"older-close, because support functions have controlled
Lipschitz norm.  Its ordinary area density $s_j(h)=gh^{p-1}$ is then close
to 1.  The classical $p=1$ Christoffel--Minkowski problem recovers $h$ with
two derivatives, modulo translation; the $C^0$ hypothesis controls that
translation.  This places the solution inside the full inverse-function
neighborhood.

\section{Non-even local theorem}

Let $\cH^{2,\alpha}$ be the open subset of $C^{2,\alpha}(\Sn)$ consisting
of positive strict support functions, and define
\[
  \Xi_p:\cH^{2,\alpha}\longrightarrow C^\alpha(\Sn),
  \qquad \Xi_p(h)=h^{1-p}s_j(h,\cdot).
\]

\begin{theorem}[full local branch and $C^0$-local uniqueness]\label{thm:main}
Fix $n\ge3$, $1\le j\le n-2$, $0\le p<1$, and $0<\alpha<1$.
Then the following hold.
\begin{enumerate}
\item There are neighborhoods $U$ of $1$ in $C^{2,\alpha}(\Sn)$ and
  $V$ of $1$ in $C^\alpha(\Sn)$ such that
  $\Xi_p:U\to V$ is a $C^1$ diffeomorphism.  No parity condition is imposed.
\item After shrinking $V$, there is $\varepsilon>0$ such that the branch
  solution $h_g\in U$ is the only classical support-function solution of
  \eqref{eq:PDE} satisfying $\|h-1\|_{C^0}<\varepsilon$.
\item If $g\in V$ is even, then $h_g$ is even.
\end{enumerate}
These statements also hold at $(p,j)=(0,1)$.
\end{theorem}

\begin{proof}
The derivative at the identity of the normalized $j$th elementary symmetric
function on $(n-1)\times(n-1)$ matrices is $j/(n-1)$ times the trace.
Consequently
\begin{equation}\label{eq:lin}
 D\Xi_p(1)[w]
 =\frac{j}{n-1}\Delta_{\Sn}w+(j+1-p)w.
\end{equation}
On spherical harmonics of degree $k$, where
$-\Delta_{\Sn}$ has eigenvalue $k(k+n-2)$, the multiplier in
\eqref{eq:lin} is
\begin{equation}\label{eq:mu}
 \mu_k=j+1-p-\frac{j}{n-1}k(k+n-2).
\end{equation}
We have $\mu_0=j+1-p>0$ and $\mu_1=1-p>0$.  For $k\ge2$,
\[
 \mu_k\le j+1-p-\frac{2jn}{n-1}<0,
\]
because
$(n-1)(j+1)-2jn=n-1-j(n+1)\le-2$.
Thus \eqref{eq:lin} has no kernel.  Schauder theory makes it a bounded
isomorphism $C^{2,\alpha}\to C^\alpha$, and the Banach inverse function
theorem proves part 1.  Notice that the calculation includes $(p,j)=(0,1)$.

We prove part 2.  Fix $0<\beta<\alpha$.  If support functions $h$ are
uniformly $C^0$-close to 1, their outer radii, hence their Lipschitz norms,
are uniformly bounded.  The elementary interpolation estimate
\begin{equation}\label{eq:interp}
 [h-1]_{C^\beta}
 \le C\|h-1\|_{C^0}^{1-\beta}
       \operatorname{Lip}(h-1)^\beta
\end{equation}
shows that $h$ is $C^\beta$-close to 1.  If also $g$ is $C^\alpha$-close
to 1, then
\begin{equation}\label{eq:fclose}
 f:=s_j(h,\cdot)=g h^{p-1}
\end{equation}
is $C^\beta$-close to 1.

Consider now the classical area-measure map $\Psi(h)=s_j(h,\cdot)$.
Its derivative at 1 is
\begin{equation}\label{eq:p1lin}
 D\Psi(1)[w]=\frac{j}{n-1}(\Delta_{\Sn}+n-1)w.
\end{equation}
The kernel of \eqref{eq:p1lin} consists exactly of degree-one harmonics,
which are translations $u\mapsto a\cdot u$.  Restrict the domain to
functions with zero degree-one component and the target to functions with
zero first moment.  (Every $j$th area measure has zero first moment.)
Then \eqref{eq:p1lin} is an isomorphism.  A second application of the
inverse function theorem to \eqref{eq:fclose} supplies a centered support
function $\widehat h$, $C^{2,\beta}$-close to 1, with
$s_j(\widehat h)=f$.

Classical uniqueness for the $j$th area-measure problem says that two
bodies with the same $S_j$ differ by a translation \cite{Schneider}.
Hence
\[
 h(u)=\widehat h(u)+a\cdot u.
\]
The degree-one component of a function $C^0$-close to 1 is small, so
$|a|$ is small.  Choosing the $C^0$ and data neighborhoods sufficiently
small places $h$ in the inverse-function neighborhood obtained from part 1
with exponent $\beta$.  The $C^{2,\alpha}$ branch solution also belongs to
that neighborhood, so its $C^{2,\beta}$ uniqueness yields $h=h_g$.

Finally, if $g$ is even then $u\mapsto h_g(-u)$ is another solution in
$U$.  Part 1 forces $h_g(-u)=h_g(u)$, proving part 3.
\end{proof}

\section{What remains global}

Theorem \ref{thm:main} removes evenness from the entire near-ball branch,
but it does not show that every solution for small data lies on that branch.
The source obtains this global conclusion in the symmetric class by a
$C^0$ non-collapsing estimate and compactness.  Its upper-bound argument
uses $h(u)\ge R|u\cdot v|$ and an origin-centered John ellipsoid, both
consequences of central symmetry.  Without symmetry one has only a
one-sided support bound, and the negative-power density $h^{p-1}$ can
concentrate on the opposite hemisphere.

The centroid identity
\[
  \int_{\Sn}u\,g(u)h(u)^{p-1}\,du=0
\]
does determine the translation uniquely inside a fixed translation class:
its derivative with respect to the translation vector is negative definite.
It does not prevent degeneration of the centered shape.  Adapting the
constant-data maximum-principle proof of Chen \cite{Chen} introduces first
and second derivatives of $g$ and likewise does not presently give the
scale-free $C^\alpha$ compactness estimate.  This is the precise remaining
obstruction to the source's full non-even question.

\section{Novelty and review status}

A bounded search used the exact title and question, together with
non-symmetric, near-isotropic, and $L_p$--Christoffel--Minkowski uniqueness
terms.  It found the source and the constant-data theorem, but no later
answer or explicit statement of Theorem \ref{thm:main}.  Novelty confidence
is modest: the full-space inverse-function step is short and may be regarded
as implicit by specialists.  The $C^0$-local upgrade and the exact isolation
of the global obstruction are the substantive additions.

\medskip
\noindent\textbf{Human-review focus.}
Check the normalization in \eqref{eq:lin}; the zero-first-moment target for
the $p=1$ map; and the use of uniqueness of $S_j$ up to translation.
Do not read this packet as a global uniqueness theorem.

\begin{thebibliography}{9}
\bibitem{PS}
K.~Patsalos and C.~Saroglou,
\emph{A note on the $L_p$--Brunn--Minkowski inequality for intrinsic volumes
and the $L_p$--Christoffel--Minkowski problem}, arXiv:2411.17896.

\bibitem{Chen}
L.~Chen,
\emph{Uniqueness of solutions to $L_p$--Christoffel--Minkowski problem for
$p<1$}, arXiv:1905.11043.

\bibitem{Schneider}
R.~Schneider,
\emph{Convex Bodies: The Brunn--Minkowski Theory}, 2nd ed.,
Cambridge University Press, 2014.
\end{thebibliography}

\end{document}
